\NSPage{025}%
\begin{EESegment}{EE-TGT-P025-001}
determine exactly one \NSi{NS-F-P025-001}{q>|z|^{1/D}}. Here is why. At that
lower endpoint, \NSi{NS-F-P025-002}{q-z^2q^{2h}} is zero. It tends to
infinity, and throughout this interval its derivative is
\NSi{NS-F-P025-003}{1-2hz^2q^{2h-1}=L\geq 1-2h>0}. This also gives
\NSi{NS-F-P025-004}{|\eta|<1}. The endpoint values
\NSi{NS-F-P025-005}{\eta=\pm1} are understood as one-sided limits.
\end{EESegment}

\begin{EESegment}{EE-TGT-P025-002}
Throughout the profile construction,
\NSi{NS-F-P025-006}{\mathcal D_Xf=X\partial_Xf} means differentiation in the
logarithmic radial coordinate, with \NSi{NS-F-P025-007}{\eta} held fixed. The
next step in computing the residual of the field from
Section~\ref{sec:3.1} is to find the operators that turn time and axial
derivatives into derivatives in \NSi{NS-F-P025-008}{(X,\eta)}.
\end{EESegment}

\begin{lemma}\label{lem:4.1}
\begin{EESegment}{EE-TGT-P025-003}
For any smooth profile \NSi{NS-F-P025-009}{f} and any real number
\NSi{NS-F-P025-010}{b},
\end{EESegment}
\NSFormula{NS-F-P025-011}{display}{4.2}
\begin{equation}
\begin{gathered}
 \partial_t(q^bf)=q^{b-1}T_bf,\qquad
 \partial_z(q^bf)=q^{b-D}Z_bf,\qquad
 \left\{
 \begin{aligned}
 T_bf&=L^{-1}(-bf+D\eta f_\eta+\mathcal D_Xf),\\
 Z_bf&=L^{-1}(2b\eta f+df_\eta-2\eta \mathcal D_Xf).
 \end{aligned}
 \right.
\end{gathered}
\tag{4.2}\label{eq:4.2}
\end{equation}
\end{lemma}

\begin{proof}
\begin{EESegment}{EE-TGT-P025-004}
Differentiate the equation for \NSi{NS-F-P025-012}{q}, then differentiate
\NSi{NS-F-P025-013}{\eta=zq^{-D}} and
\NSi{NS-F-P025-014}{X=s/q}. This gives
\end{EESegment}
\NSFormula{NS-F-P025-015}{display}{none}
\[
\begin{aligned}
 q_t&=-L^{-1},& \eta_t&=D\eta/(qL),& X_t&=X/(qL),\\
 q_z&=2\eta q^{1-D}/L,& \eta_z&=d/(q^DL),& X_z&=-2\eta X/(q^DL).
\end{aligned}
\]
\begin{EESegment}{EE-TGT-P025-005}
The middle entry in the second row uses
\NSi{NS-F-P025-016}{L-2D\eta^2=d}. The formulas in the lemma now follow from
the chain rule.
\end{EESegment}
\end{proof}

\begin{EESegment}{EE-TGT-P025-006}
The leading field from Section~\ref{sec:3.1} must be smooth across the axis.
Factor out the linear radial vanishing of its azimuthal component, choose a fixed amplitude normalization
\NSi{NS-F-P025-017}{\mathsf C>1} and a smooth scalar
\NSi{NS-F-P025-018}{\phi}, and write
\end{EESegment}
\NSFormula{NS-F-P025-019}{display}{4.3}
\begin{equation}
 u_\theta^{(0)}=q^{-A}E,\qquad
 u_z^{(0)}=q^{-A}U,\qquad
 ru_r^{(0)}=V_0,\qquad
 p^{(0)}=q^{-2A}\Pi,\qquad
 E=\mathsf C^{-1}\sqrt{2X}\,\phi.
\tag{4.3}\label{eq:4.3}
\end{equation}
\begin{EESegment}{EE-TGT-P025-007}
The proof of Theorem~\ref{thm:4.6} chooses the normalization
\NSi{NS-F-P025-020}{\mathsf C}. As in Section~\ref{sec:3.1},
\NSi{NS-F-P025-021}{E,U,V_0,\Pi} are functions of
\NSi{NS-F-P025-022}{(X,\eta)}. The factor
\NSi{NS-F-P025-023}{\sqrt{2X}} makes the azimuthal velocity vanish at the
axis, while the smoothness condition is put on
\NSi{NS-F-P025-024}{\phi}. For the profiles built here, the following
conditions give regularity at the axis in the sense of
Definition~\ref{def:3.2}.
\end{EESegment}
\NSFormula{NS-F-P025-025}{display}{4.4}
\begin{equation}
 E=\sqrt{2X}\,F,\qquad V_0=Xv_0,\qquad
 F=\phi/\mathsf C,\qquad
 U,v_0,\Pi\in C^\infty([0,X_c]\times[-1,1])
\tag{4.4}\label{eq:4.4}
\end{equation}
\begin{EESegment}{EE-TGT-P025-008}
Here \NSi{NS-F-P025-026}{X_c>0}. Smoothness on this closed rectangle means
that every mixed \NSi{NS-F-P025-027}{X,\eta} derivative extends continuously
to the boundary from the one side that lies inside the rectangle. To see why
this gives Cartesian smoothness, write
\NSi{NS-F-P025-028}{x_1=r\cos\theta,\ x_2=r\sin\theta}. Then
\NSi{NS-F-P025-029}{X=(x_1^2+x_2^2)/(2q)}, and
\end{EESegment}
\NSFormula{NS-F-P025-030}{display}{4.5}
\begin{equation}
\begin{aligned}
 u_1^{(0)}&=\frac{v_0}{2q}x_1-q^{-A-1/2}Fx_2,
 &u_2^{(0)}&=\frac{v_0}{2q}x_2+q^{-A-1/2}Fx_1,\\
 u_3^{(0)}&=q^{-A}U,
 &p^{(0)}&=q^{-2A}\Pi.
\end{aligned}
\tag{4.5}\label{eq:4.5}
\end{equation}
\begin{EESegment}{EE-TGT-P025-009}
Every profile in this display is evaluated at
\NSi{NS-F-P025-031}{(X,\eta)}. The quantities
\NSi{NS-F-P025-032}{q,\eta} are smooth functions of
\NSi{NS-F-P025-033}{(t,z)}, and \NSi{NS-F-P025-034}{q>0} whenever
\NSi{NS-F-P025-035}{t<1}. So the coefficients in \eqref{eq:4.5} are
smooth functions of \NSi{NS-F-P025-036}{(x_1,x_2,z,t)}. Notice in particular
that \NSi{NS-F-P025-037}{E} contains a
\NSi{NS-F-P025-038}{\sqrt X} factor. At
\NSi{NS-F-P025-039}{X=0}, smoothness is imposed on
\NSi{NS-F-P025-040}{F}.
\end{EESegment}

\begin{EESegment}{EE-TGT-P025-010}
Choose \NSi{NS-F-P025-041}{E,U} first. Incompressibility then determines
\NSi{NS-F-P025-042}{V_0}, and the leading radial momentum balance determines
\NSi{NS-F-P025-043}{\Pi} once its value on the axis has been fixed. The
notation paragraph at the end of Section~\ref{sec:3} defined the radial
average of a smooth radial scalar. Its smooth extension to the axis is
specified by
\end{EESegment}
\NSFormula{NS-F-P025-044}{display}{4.6}
\begin{equation}
 \mathcal A_X(f)(X,\eta)=\frac1X\int_0^X f(x,\eta)\,dx,
 \qquad
 \mathcal A_X(f)(0,\eta)=f(0,\eta).
\tag{4.6}\label{eq:4.6}
\end{equation}

\NSPage{026}%
\begin{EESegment}{EE-TGT-P026-001}
The incompressibility and centrifugal-pressure balances from
Section~\ref{sec:3.1} now give explicit formulas for the profiles.
\end{EESegment}
\NSFormula{NS-F-P026-001}{display}{4.7}
\begin{equation}
 V_0=\frac{X}{L}\bigl(2\eta U-2D\eta\mathcal A_X(U)
      -d\partial_\eta\mathcal A_X(U)\bigr),
 \qquad
 \Pi_X=\frac{E^2}{2X}.
\tag{4.7}\label{eq:4.7}
\end{equation}
\begin{EESegment}{EE-TGT-P026-002}
For the first formula, put \NSi{NS-F-P026-002}{A+D=1} into
\NSi{NS-F-P026-003}{\partial_s(ru_r^{(0)})+\partial_z u_z^{(0)}=0}. The
result is
\NSi{NS-F-P026-004}{(V_0)_X=L^{-1}(2A\eta U-dU_\eta+2\eta XU_X)}. Integrate
it and set the value on the axis to zero to get the displayed formula.
\end{EESegment}

\begin{EESegment}{EE-TGT-P026-003}
For the second formula, the coefficients of
\NSi{NS-F-P026-005}{\partial_rp^{(0)}} and
\NSi{NS-F-P026-006}{(u_\theta^{(0)})^2/r} are
\NSi{NS-F-P026-007}{\sqrt{2X}\Pi_X} and
\NSi{NS-F-P026-008}{E^2/\sqrt{2X}}, respectively. Both carry the same power
\NSi{NS-F-P026-009}{q^{-2A-1/2}}. The other radial terms carry at least one
extra factor \NSi{NS-F-P026-010}{q^{2h}}. The same is true of axial
viscosity when it is compared with radial viscosity in the tangential
equations. This is why the background corrections in Section~\ref{sec:3}
use the powers \NSi{NS-F-P026-011}{q^{2nh}}. These comparisons are being
made with the profile coordinates fixed. Bounds for physical derivatives of
the finished fields are proved separately. The full residual still includes
the higher-order terms just named, and Section~\ref{sec:5} deals with them.
\end{EESegment}

\begin{EESegment}{EE-TGT-P026-004}
After imposing incompressibility and the leading radial balance, the
tangential residual still has to be written as a radial stress divergence,
as Section~\ref{sec:3} requires. Start by integrating its inviscid terms and
call the resulting scalar profiles \NSi{NS-F-P026-012}{Q_s,N_s}. Then add
the part that comes from radial viscosity. The definitions below use only
\NSi{NS-F-P026-013}{E,U,\Pi}; Proposition~\ref{prop:4.2} checks that they
give the required stress identity.
\end{EESegment}

\begin{EESegment}{EE-TGT-P026-005}
Take profiles that satisfy \eqref{eq:4.7}, with
\NSi{NS-F-P026-014}{\phi,U,\Pi} smooth at the axis and
\NSi{NS-F-P026-015}{\phi>0}, and set
\end{EESegment}
\NSFormula{NS-F-P026-016}{display}{4.8}
\begin{equation}
\begin{gathered}
 H=\sqrt{2X}E,\qquad F=\frac{E}{\sqrt{2X}},\qquad l=\mathcal D_X\log H,\\
 W=1-2D\eta\mathcal A_X(U)-d\partial_\eta\mathcal A_X(U),
 \qquad H_c=D\eta+dU.
\end{gathered}
\tag{4.8}\label{eq:4.8}
\end{equation}
\begin{EESegment}{EE-TGT-P026-006}
Here \NSi{NS-F-P026-017}{H} is the profile of the angular momentum
\NSi{NS-F-P026-018}{ru_\theta^{(0)}}. The quantity
\NSi{NS-F-P026-019}{F} is what remains after taking the factor
\NSi{NS-F-P026-020}{\sqrt{2X}} out of
\NSi{NS-F-P026-021}{E}, and \NSi{NS-F-P026-022}{l} is the logarithmic radial
slope of \NSi{NS-F-P026-023}{H}. The two coefficients
\NSi{NS-F-P026-024}{W,H_c} gather the transport terms that appear in
\eqref{eq:4.14}.
\end{EESegment}

\begin{EESegment}{EE-TGT-P026-007}
Define \NSi{NS-F-P026-025}{Q_s,N_s} by the radial equations below, choosing
the solutions that extend smoothly to \NSi{NS-F-P026-026}{X=0}.
\end{EESegment}
\NSFormula{NS-F-P026-027}{display}{4.9}
\begin{equation}
\begin{aligned}
 \mathcal D_XQ_s+(1+l)Q_s&=S_q,& \mathcal D_XN_s+N_s&=S_n,\\
 S_q&=-Wl-h(1-2\eta U)-H_c(\log E)_\eta,\\
 S_n&=-W\mathcal D_XU-A(1-2\eta U)U-H_cU_\eta
      -d\Pi_\eta+4A\eta\Pi+2\eta \mathcal D_X\Pi.
\end{aligned}
\tag{4.9}\label{eq:4.9}
\end{equation}
\begin{EESegment}{EE-TGT-P026-008}
These are first-order equations in the radial variable. After each equation is divided by that
variable, multiplying the first by \NSi{NS-F-P026-028}{XH} turns its left side into the radial
derivative of the product of that factor with the first unknown. Multiplying the second by
\NSi{NS-F-P026-029}{X} does the same for the second unknown. That is why these expressions are
called integrating factors. Before smoothness at the
axis is imposed, the constants of integration may depend on
\NSi{NS-F-P026-030}{\eta}, so the solutions have the form
\end{EESegment}
\NSFormula{NS-F-P026-031}{display}{4.10}
\begin{equation}
 Q_s=\frac{C_Q(\eta)+\int_0^X H(x,\eta)S_q(x,\eta)\,dx}{XH(X,\eta)},
 \qquad
 N_s=\frac{C_N(\eta)+\int_0^X S_n(x,\eta)\,dx}{X}.
\tag{4.10}\label{eq:4.10}
\end{equation}
\begin{EESegment}{EE-TGT-P026-009}
Smooth extension to the axis forces
\NSi{NS-F-P026-032}{C_Q=C_N=0}. The reason is that
\NSi{NS-F-P026-033}{H=2X\phi/\mathsf C} near the axis. The first constant,
if nonzero, would give an \NSi{NS-F-P026-034}{X^{-2}} singularity in
\NSi{NS-F-P026-035}{Q_s}, while a nonzero second constant would give an
\NSi{NS-F-P026-036}{X^{-1}} singularity in \NSi{NS-F-P026-037}{N_s}.
With both constants set to zero, the change of
variables \NSi{NS-F-P026-038}{x=sX} gives
\end{EESegment}
\NSFormula{NS-F-P026-039}{display}{none}
\[
 Q_s=\frac{1}{F(X,\eta)}\int_0^1 sF(sX,\eta)S_q(sX,\eta)\,ds,
 \qquad
 N_s=\int_0^1 S_n(sX,\eta)\,ds.
\]
\begin{EESegment}{EE-TGT-P026-010}
The sources are smooth at \NSi{NS-F-P026-040}{X=0}, because
\NSi{NS-F-P026-041}{l=1+XF_X/F} and
\NSi{NS-F-P026-042}{(\log E)_\eta=F_\eta/F}, so the formulas extend
smoothly to the axis with
\NSi{NS-F-P026-043}{Q_s(0,\eta)=S_q(0,\eta)/2} and
\NSi{NS-F-P026-044}{N_s(0,\eta)=S_n(0,\eta)}.
\end{EESegment}

\NSPage{027}%
\begin{EESegment}{EE-TGT-P027-001}
Put the two stress coefficients in the order
\NSi{NS-F-P027-001}{(r\theta,rz)}, and define them by
\end{EESegment}
\NSFormula{NS-F-P027-002}{display}{4.11}
\begin{equation}
\begin{gathered}
 \mathcal T_0=F(p_s-\mathfrak{s}),
 \qquad
 p_s=\left(\frac{XQ_s}{L},\frac{XN_s}{LE}\right),\\
 \mathfrak{s}=(a,-b_s),
 \qquad
 a=1-2\mathcal D_X\log E=2-2l,
 \qquad
 b_s=\frac{2\mathcal D_XU}{E}.
\end{gathered}
\tag{4.11}\label{eq:4.11}
\end{equation}
\begin{EESegment}{EE-TGT-P027-002}
The symbol \NSi{NS-F-P027-003}{p_s} now means the two-component vector of
integrated inviscid terms. It is not the pressure
\NSi{NS-F-P027-004}{p^{(0)}}. The vector
\NSi{NS-F-P027-005}{\mathfrak{s}=(a,-b_s)} records the cylindrical radial
shear, with its normalization fixed by
\end{EESegment}
\NSFormula{NS-F-P027-006}{display}{none}
\[
 \left(\partial_ru_\theta^{(0)}-\frac{u_\theta^{(0)}}{r},
       \partial_ru_z^{(0)}\right)
 =-q^{-A-1/2}F\mathfrak{s}.
\]
\begin{EESegment}{EE-TGT-P027-003}
In \eqref{eq:4.11}, \NSi{NS-F-P027-007}{Fp_s} is the inviscid part and
\NSi{NS-F-P027-008}{-F\mathfrak{s}} is the part from radial viscosity. The
calculation below shows both facts. In the notation of Section~\ref{sec:3},
the formula to prove for the physical stress is
\NSi{NS-F-P027-009}{T=q^{-A-1/2}\mathcal T_0}. Its components are the
\NSi{NS-F-P027-010}{r\theta} and \NSi{NS-F-P027-011}{rz} entries needed from
the wave covariance, which is the tensor of averaged quadratic velocity
products built in Proposition~\ref{prop:7.5}.
\end{EESegment}

\begin{EESegment}{EE-TGT-P027-004}
To check the stress identity from Section~\ref{sec:3}, separate the leading
tangential terms from axial viscosity. Define the
\emph{leading tangential residuals} by
\end{EESegment}
\NSFormula{NS-F-P027-012}{display}{4.12}
\begin{equation}
 R_j^{(0)}=\mathscr R(u^{(0)},p^{(0)})\mathbin{\cdot}e_j
            +\partial_z^2u_j^{(0)},
 \qquad j\in\{\theta,z\},
\tag{4.12}\label{eq:4.12}
\end{equation}
\begin{EESegment}{EE-TGT-P027-005}
where \NSi{NS-F-P027-013}{e_\theta,e_z} are the cylindrical unit vectors and
\NSi{NS-F-P027-014}{\mathscr R} is the full momentum residual defined in the
introduction. Adding \NSi{NS-F-P027-015}{\partial_z^2u_j^{(0)}} takes the
axial-viscosity term out of the residual but leaves radial viscosity in it.
\end{EESegment}

\begin{proposition}\label{prop:4.2}
\begin{EESegment}{EE-TGT-P027-006}
The residuals \NSi{NS-F-P027-016}{R_\theta^{(0)},R_z^{(0)}} from
\eqref{eq:4.12} are the negatives of the cylindrical divergences
\NSi{NS-F-P027-017}{\partial_r+2/r} and
\NSi{NS-F-P027-018}{\partial_r+1/r}, applied to
\NSi{NS-F-P027-019}{q^{-A-1/2}\mathcal T_0}. In particular, the equations
that set the leading residual stress to zero are
\end{EESegment}
\NSFormula{NS-F-P027-020}{display}{4.13}
\begin{equation}
 -2L\frac{X\phi_{XX}+2\phi_X}{\phi}=S_q,
 \qquad
 -2L(XU_{XX}+U_X)=S_n.
\tag{4.13}\label{eq:4.13}
\end{equation}
\end{proposition}

\begin{proof}
\begin{EESegment}{EE-TGT-P027-007}
The angular equation takes its simplest form when written for the angular
momentum \NSi{NS-F-P027-021}{ru_\theta^{(0)}=q^{-h}H}.
Equations~\eqref{eq:4.2} and~\eqref{eq:4.7} give
\end{EESegment}
\NSFormula{NS-F-P027-022}{display}{4.14}
\begin{equation}
\begin{split}
 &(\partial_t+u_r^{(0)}\partial_r+u_z^{(0)}\partial_z)(q^{-h}H)\\
 &\qquad=\frac{q^{-h-1}}{L}
   \{W\mathcal D_XH+H_cH_\eta+h(1-2\eta U)H\}
   =-\frac{q^{-h-1}}{L}HS_q.
\end{split}
\tag{4.14}\label{eq:4.14}
\end{equation}
\begin{EESegment}{EE-TGT-P027-008}
The same calculation gives
\NSi{NS-F-P027-023}{-q^{-A-1}S_n/L} for the axial material derivative plus
the pressure gradient. The integrating factors for the two
radial divergences are \NSi{NS-F-P027-024}{r^2} and
\NSi{NS-F-P027-025}{r}. So integrating the negative inviscid residual outward
from the axis gives the coefficients
\NSi{NS-F-P027-026}{FXQ_s/L} and
\NSi{NS-F-P027-027}{FXN_s/(LE)}. Radial viscosity contributes
\end{EESegment}
\NSFormula{NS-F-P027-028}{display}{none}
\[
 \partial_ru_\theta^{(0)}-u_\theta^{(0)}/r=-q^{-A-1/2}Fa,
 \qquad
 \partial_ru_z^{(0)}=q^{-A-1/2}Fb_s.
\]
\begin{EESegment}{EE-TGT-P027-009}
Add these viscous terms to the radially integrated inviscid residual. This
gives \eqref{eq:4.11}, including its sign. If the stress is zero, then
\NSi{NS-F-P027-029}{XQ_s/L=a,\ N_s/L=-2U_X}. Put these two relations into
\eqref{eq:4.9} to get \eqref{eq:4.13}. On the axis,
\NSi{NS-F-P027-030}{l=1+X\phi_X/\phi} and
\NSi{NS-F-P027-031}{(\log E)_\eta=\phi_\eta/\phi} are smooth, so
\NSi{NS-F-P027-032}{Q_s,N_s} are regular on the axis in the scalar-profile sense
of Definition~\ref{def:3.2}.
\end{EESegment}
\end{proof}

\NSPage{028}%
\begin{EESegment}{EE-TGT-P028-001}
\subsection{The cumulative radial integrals}\label{sec:4.2}

The stress profile \NSi{NS-F-P028-001}{\mathcal T_0} in \eqref{eq:4.11}
depends, through the integrals in \eqref{eq:4.10}, on the profiles
\NSi{NS-F-P028-002}{E,U} at every smaller radius. This matters when an inner
profile is joined to a prescribed outer profile, as in Section~\ref{sec:3}.
The accumulated integrals must match along with the profile values. Five
integrals are enough to keep the outer pressure, radial velocity, and stress
unchanged, as Lemma~\ref{lem:4.4} shows. Integrating \eqref{eq:4.9} writes
\NSi{NS-F-P028-003}{Q_s,N_s} in terms of those five quantities.
\end{EESegment}
\NSFormula{NS-F-P028-004}{display}{4.15}
\begin{equation}
\begin{gathered}
 M=\int_0^X U\,dx,\qquad
 I=\int_0^X H\,dx,\qquad
 J=\int_0^X UH\,dx,\\
 S=\int_0^X(U^2-E^2/2)\,dx,\qquad
 C_p=\int_0^X\frac{E^2}{2x}\,dx,\qquad
 \Pi=\Pi(0,\eta)+C_p.
\end{gathered}
\tag{4.15}\label{eq:4.15}
\end{equation}
\begin{EESegment}{EE-TGT-P028-002}
Each integral depends on \NSi{NS-F-P028-005}{(X,\eta)}. Once the matching
radius is fixed, each one is a scalar function of
\NSi{NS-F-P028-006}{\eta}. The quantity \NSi{NS-F-P028-007}{C_p} is the
pressure increment from the axis to \NSi{NS-F-P028-008}{X}, while
\NSi{NS-F-P028-009}{\Pi(0,\eta)} is the prescribed pressure value at
\NSi{NS-F-P028-010}{X=0}.
\end{EESegment}

\begin{lemma}\label{lem:4.3}
\begin{EESegment}{EE-TGT-P028-003}
For the solutions of \eqref{eq:4.9} with
\NSi{NS-F-P028-011}{C_Q=C_N=0},
\end{EESegment}
\NSFormula{NS-F-P028-012}{display}{4.16}
\begin{equation}
\begin{aligned}
 Q_s&=-W+\frac{(1-h)I-D\eta I_\eta-dJ_\eta+2(h-D)\eta J}{XH},\\
 N_s&=-WU+\frac{D(M-\eta M_\eta)+4h\eta S-dS_\eta}{X}
             +4A\eta\Pi-d\Pi_\eta.
\end{aligned}
\tag{4.16}\label{eq:4.16}
\end{equation}
\end{lemma}

\begin{proof}
\begin{EESegment}{EE-TGT-P028-004}
Use \NSi{NS-F-P028-013}{\partial_X(XW)=1-2D\eta U-dU_\eta} to integrate
\NSi{NS-F-P028-014}{-W\mathcal D_XH} and
\NSi{NS-F-P028-015}{-W\mathcal D_XU}. In the angular terms,
\NSi{NS-F-P028-016}{-dU_\eta H-dUH_\eta} becomes
\NSi{NS-F-P028-017}{-d(UH)_\eta}. For the pressure term, integrate by parts
in \NSi{NS-F-P028-018}{\Pi_X=E^2/(2X)} to get
\NSi{NS-F-P028-019}{\int_0^X\Pi\,dx=X\Pi-\int_0^XE^2/2\,dx}. The source
integrals are then
\end{EESegment}
\NSFormula{NS-F-P028-020}{display}{none}
\[
\begin{aligned}
 \int_0^XHS_q\,dx
   &=-XWH+(1-h)I-D\eta I_\eta-dJ_\eta+2(h-D)\eta J,\\
 \int_0^XS_n\,dx
   &=-XWU+D(M-\eta M_\eta)+4h\eta S-dS_\eta
      +X(4A\eta\Pi-d\Pi_\eta).
\end{aligned}
\]
\begin{EESegment}{EE-TGT-P028-005}
Divide the first line by \NSi{NS-F-P028-021}{XH} and the second by
\NSi{NS-F-P028-022}{X}. This gives the two formulas in the lemma.
\end{EESegment}
\end{proof}

\begin{EESegment}{EE-TGT-P028-006}
Suppose two profiles agree beyond a joining radius. If their five cumulative
integrals also match at that radius, then all quantities obtained by radial
integration stay the same beyond it. Their common pressure value on the axis
is held fixed. If the profiles, the integrals, and their parameter
derivatives change only a little, then
\NSi{NS-F-P028-023}{Q_s,N_s,p_s} also change only a little. This estimate
does not require the radial derivatives of the two profiles to be close. The
exact identity is used to join profiles, and the estimate is used to change
their shear.
\end{EESegment}

\begin{lemma}\label{lem:4.4}
\begin{EESegment}{EE-TGT-P028-007}
Fix \NSi{NS-F-P028-024}{h\in(0,1/2)}. For
\NSi{NS-F-P028-025}{i=1,2}, let \NSi{NS-F-P028-026}{(U_i,E_i)} be smooth on
\NSi{NS-F-P028-027}{(0,\infty)\times[-1,1]}, with
\NSi{NS-F-P028-028}{E_i>0}. Suppose the five integrals in
\eqref{eq:4.15} are smooth functions there, including their one-sided
\NSi{NS-F-P028-029}{\eta}-derivatives at
\NSi{NS-F-P028-030}{\eta=\pm1}. Write
\end{EESegment}
\NSFormula{NS-F-P028-031}{display}{none}
\[
 H_i=\sqrt{2X}E_i,\qquad
 \mathfrak m_i=(M_i,I_i,J_i,S_i,C_{p,i}).
\]
\begin{EESegment}{EE-TGT-P028-008}
Fix one smooth function \NSi{NS-F-P028-032}{\Pi_{\mathrm{ax}}(\eta)} on
\NSi{NS-F-P028-033}{[-1,1]}, and set
\end{EESegment}
\NSFormula{NS-F-P028-034}{display}{none}
\[
 \Pi_i(X,\eta)=\Pi_{\mathrm{ax}}(\eta)+C_{p,i}(X,\eta),
 \qquad
 \Pi_i(0,\eta)=\Pi_{\mathrm{ax}}(\eta),
 \qquad i=1,2.
\]
\begin{EESegment}{EE-TGT-P028-009}
This gives both pressures the same constant of integration on the axis. The
velocity profiles are not required to agree near the axis. Define
\NSi{NS-F-P028-035}{V_{0,i}} by \eqref{eq:4.7},
\NSi{NS-F-P028-036}{Q_{s,i},N_{s,i}} by \eqref{eq:4.16}, and
\NSi{NS-F-P028-037}{p_{s,i},a_i,b_{s,i},\mathcal T_{0,i}} by
\eqref{eq:4.11}. The formulas \eqref{eq:4.16} used here are the ones already
derived with \NSi{NS-F-P028-038}{C_Q=C_N=0}. For any one of these quantities
\NSi{NS-F-P028-039}{f}, write
\NSi{NS-F-P028-040}{\Delta f=f_2-f_1}. There are two conclusions.
\end{EESegment}
\NSPage{029}%
\begin{enumerate}[label=\textup{(\roman*)},leftmargin=2.5em]
\item
\begin{EESegment}{EE-TGT-P029-001}
Suppose that for some \NSi{NS-F-P029-001}{X_h>0},
\end{EESegment}
\NSFormula{NS-F-P029-002}{display}{none}
\[
 (U_1,E_1)=(U_2,E_2)\quad\text{for }X\geq X_h,
 \qquad
 \Delta\mathfrak m(X_h,\eta)=0\quad(-1\leq\eta\leq1).
\]
\begin{EESegment}{EE-TGT-P029-002}
Then \NSi{NS-F-P029-003}{\mathfrak m_1=\mathfrak m_2}. All the corresponding
quantities
\NSi{NS-F-P029-004}{\Pi_i,V_{0,i},Q_{s,i},N_{s,i},p_{s,i},a_i,b_{s,i},
\mathcal T_{0,i}} also agree whenever \NSi{NS-F-P029-005}{X\geq X_h} and
\NSi{NS-F-P029-006}{\eta\in[-1,1]}.
\end{EESegment}

\item
\begin{EESegment}{EE-TGT-P029-003}
Fix \NSi{NS-F-P029-007}{0<X_0<X_1<\infty}, an integer
\NSi{NS-F-P029-008}{k\geq0}, and \NSi{NS-F-P029-009}{e_{\min}>0}. On the
rectangle \NSi{NS-F-P029-010}{\mathcal R=[X_0,X_1]\times[-1,1]}, put
\end{EESegment}
\NSFormula{NS-F-P029-011}{display}{none}
\[
 \|f\|_j:=\max_{0\leq r\leq j}
           \sup_{(X,\eta)\in\mathcal R}|\partial_\eta^rf(X,\eta)|,
 \qquad j\geq0,
\]
\begin{EESegment}{EE-TGT-P029-004}
and for a tuple use the largest of its component norms. If
\NSi{NS-F-P029-012}{E_i\geq e_{\min}} on
\NSi{NS-F-P029-013}{\mathcal R}, then
\end{EESegment}
\NSFormula{NS-F-P029-014}{display}{4.17}
\begin{equation}
 \|\Delta(Q_s,N_s,p_s)\|_k
 \leq C_k\bigl(\|\Delta(U,E)\|_{k+1}
                 +\|\Delta\mathfrak m\|_{k+1}\bigr).
\tag{4.17}\label{eq:4.17}
\end{equation}
\begin{EESegment}{EE-TGT-P029-005}
The constant \NSi{NS-F-P029-015}{C_k} depends only on
\NSi{NS-F-P029-016}{k,h,X_0,X_1,e_{\min}}, the
\NSi{NS-F-P029-017}{C^{k+1}}-norm of
\NSi{NS-F-P029-018}{\Pi_{\mathrm{ax}}}, and one common bound for
\NSi{NS-F-P029-019}{\|(U_i,E_i,\mathfrak m_i)\|_{k+1}},
\NSi{NS-F-P029-020}{i=1,2}. The estimate contains no radial derivative of
any difference.
\end{EESegment}
\end{enumerate}
\end{lemma}

\begin{proof}
\begin{EESegment}{EE-TGT-P029-006}
The definitions of the five cumulative integrals give
\end{EESegment}
\NSFormula{NS-F-P029-021}{display}{4.18}
\begin{equation}
\begin{aligned}
 \Delta M&=\int_0^X(U_2-U_1)\,dx,\\
 \Delta I&=\int_0^X(H_2-H_1)\,dx,\\
 \Delta J&=\int_0^X(U_2H_2-U_1H_1)\,dx,\\
 \Delta S&=\int_0^X\bigl(U_2^2-U_1^2-(E_2^2-E_1^2)/2\bigr)\,dx,\\
 \Delta C_p&=\int_0^X\frac{E_2^2-E_1^2}{2x}\,dx.
\end{aligned}
\tag{4.18}\label{eq:4.18}
\end{equation}
\begin{EESegment}{EE-TGT-P029-007}
The quantities on the left are evaluated at
\NSi{NS-F-P029-022}{(X,\eta)}, while each integrand on the right is evaluated
at \NSi{NS-F-P029-023}{(x,\eta)}. In part~\textup{(i)}, the profiles
\NSi{NS-F-P029-024}{U_i,E_i} agree for
\NSi{NS-F-P029-025}{X\geq X_h}, so every integrand in \eqref{eq:4.18}
vanishes there. It follows that
\end{EESegment}
\NSFormula{NS-F-P029-026}{display}{none}
\[
 \Delta\mathfrak m(X,\eta)=\Delta\mathfrak m(X_h,\eta)=0
 \qquad(X\geq X_h,\ -1\leq\eta\leq1).
\]
\begin{EESegment}{EE-TGT-P029-008}
This is an identity between functions of \NSi{NS-F-P029-027}{\eta}, so their
\NSi{NS-F-P029-028}{\eta}-derivatives agree too. The common pressure value
on the axis gives \NSi{NS-F-P029-029}{\Delta\Pi=\Delta C_p=0}. Then
\eqref{eq:4.7} gives \NSi{NS-F-P029-030}{\Delta V_0=0}, and \eqref{eq:4.16}
gives \NSi{NS-F-P029-031}{\Delta Q_s=\Delta N_s=0}. The profiles agree on a
radial interval, so their radial derivatives agree there as well. The
definitions in \eqref{eq:4.11} now give equality of
\NSi{NS-F-P029-032}{p_s,a,b_s,\mathcal T_0}.
\end{EESegment}

\begin{EESegment}{EE-TGT-P029-009}
For part~\textup{(ii)}, write the quantities that enter \eqref{eq:4.16} as
\end{EESegment}
\NSFormula{NS-F-P029-033}{display}{none}
\[
 H_i=\sqrt{2X}E_i,
 \qquad
 W_i=1-\frac{2D\eta M_i+d\partial_\eta M_i}{X},
 \qquad
 \Pi_i=\Pi_{\mathrm{ax}}+C_{p,i}.
\]
\begin{EESegment}{EE-TGT-P029-010}
Their differences are
\end{EESegment}
\NSFormula{NS-F-P029-034}{display}{none}
\[
 \Delta H=\sqrt{2X}\Delta E,
 \qquad
 \Delta W=-\frac{2D\eta\Delta M+d\partial_\eta\Delta M}{X},
 \qquad
 \Delta\Pi=\Delta C_p.
\]
\begin{EESegment}{EE-TGT-P029-011}
The last equality uses the common pressure value on the axis. On
\NSi{NS-F-P029-035}{\mathcal R}, every denominator stays away from zero,
because \NSi{NS-F-P029-036}{X\geq X_0},
\NSi{NS-F-P029-037}{L\geq1-2h>0}, and
\NSi{NS-F-P029-038}{H_i\geq\sqrt{2X_0}e_{\min}}. For example,
\end{EESegment}
\NSFormula{NS-F-P029-039}{display}{none}
\[
 \Delta(H^{-1})=-\frac{\Delta H}{H_1H_2},
 \qquad
 \Delta(E^{-1})=-\frac{\Delta E}{E_1E_2}.
\]
\begin{EESegment}{EE-TGT-P029-012}
Apply the product rule and these two reciprocal identities first to
\eqref{eq:4.16}, then to the formula for
\NSi{NS-F-P029-040}{p_s} in \eqref{eq:4.11}. After no more than
\NSi{NS-F-P029-041}{k} derivatives with respect to
\NSi{NS-F-P029-042}{\eta}, every term contains a difference of
\NSi{NS-F-P029-043}{U,E}, or a difference of
\end{EESegment}
\NSPage{030}%
\begin{EESegment}{EE-TGT-P030-001}
a component of \NSi{NS-F-P030-001}{\mathfrak m}, through order at most
\NSi{NS-F-P030-002}{k+1}. The remaining factors are bounded by the common
bounds stated in the lemma. This proves \eqref{eq:4.17}, including the stated
dependence of its constant. None of these formulas contains a radial
derivative of an input.
\end{EESegment}
\end{proof}

\begin{EESegment}{EE-TGT-P030-002}
Part~\textup{(ii)} is used for the radial modulation in
Proposition~\ref{prop:C.2}. The profiles \NSi{NS-F-P030-003}{U,E}, the five
integrals \NSi{NS-F-P030-004}{\mathfrak m}, and their
\NSi{NS-F-P030-005}{\eta}-derivatives stay close even when the radial shears
\NSi{NS-F-P030-006}{a,b_s} change a lot. The estimate controls
\NSi{NS-F-P030-007}{p_s}. The full stress
\NSi{NS-F-P030-008}{\mathcal T_0=F(p_s-\mathfrak{s})} also depends on those
shears.
\end{EESegment}

\begin{EESegment}{EE-TGT-P030-003}
To use this comparison at a large joining radius
\NSi{NS-F-P030-009}{\rho}, set
\NSi{NS-F-P030-010}{\widehat U(x,\eta)=U(\rho x,\eta)} and
\NSi{NS-F-P030-011}{\widehat E(x,\eta)=E(\rho x,\eta)}. Define the rest of
the hatted quantities by
\end{EESegment}
\NSFormula{NS-F-P030-012}{display}{4.19}
\begin{equation}
 X=\rho x,\qquad H=\sqrt\rho\,\widehat H,\qquad
 (M,I,J,S,C_p)
   =(\rho\widehat M,\rho^{3/2}\widehat I,\rho^{3/2}\widehat J,
     \rho\widehat S,\widehat C_p).
\tag{4.19}\label{eq:4.19}
\end{equation}
\begin{EESegment}{EE-TGT-P030-004}
Here the unhatted functions are evaluated at
\NSi{NS-F-P030-013}{(\rho x,\eta)}, while the hatted functions are evaluated
at \NSi{NS-F-P030-014}{(x,\eta)}. Put
\end{EESegment}
\NSFormula{NS-F-P030-015}{display}{none}
\[
\begin{gathered}
 \widehat{\mathfrak m}
   =(\widehat M,\widehat I,\widehat J,\widehat S,\widehat C_p),\\
 (\widehat Q_s,\widehat N_s)(x,\eta)=(Q_s,N_s)(\rho x,\eta),\\
 \widehat p_s(x,\eta)=\rho^{-1}p_s(\rho x,\eta).
\end{gathered}
\]
\begin{EESegment}{EE-TGT-P030-005}
Substitute these expressions into \eqref{eq:4.16} and \eqref{eq:4.11}.
The formulas for
\NSi{NS-F-P030-016}{\widehat Q_s,\widehat N_s,\widehat p_s} are then the
same, with \NSi{NS-F-P030-017}{X} replaced by
\NSi{NS-F-P030-018}{x}. On any fixed rectangle
\NSi{NS-F-P030-019}{[x_0,x_1]\times[-1,1]}, where
\NSi{NS-F-P030-020}{0<x_0<x_1}, this gives
\end{EESegment}
\NSFormula{NS-F-P030-021}{display}{none}
\[
 \|\Delta(\widehat Q_s,\widehat N_s,\widehat p_s)\|_k
 \leq C_k\bigl(\|\Delta(\widehat U,\widehat E)\|_{k+1}
                 +\|\Delta\widehat{\mathfrak m}\|_{k+1}\bigr),
\]
\begin{EESegment}{EE-TGT-P030-006}
where the norms are taken on this fixed
\NSi{NS-F-P030-022}{x}-rectangle. The constant does not depend on
\NSi{NS-F-P030-023}{\rho}, provided that the positive lower bound for
\NSi{NS-F-P030-024}{\widehat E_i}, the bounds for
\NSi{NS-F-P030-025}{\|(\widehat U_i,\widehat E_i,
\widehat{\mathfrak m}_i)\|_{k+1}}, and the
\NSi{NS-F-P030-026}{C^{k+1}}-bound for the common
\NSi{NS-F-P030-027}{\Pi_{\mathrm{ax}}} are all independent of
\NSi{NS-F-P030-028}{\rho}, with \NSi{NS-F-P030-029}{h} fixed.
\end{EESegment}

\begin{EESegment}{EE-TGT-P030-007}
The moment corrections used below are finite linear combinations of smooth
functions supported on short radial intervals. Their changes to
\NSi{NS-F-P030-030}{\mathfrak m} depend linearly or quadratically on the
coefficients. Lemmas~\ref{lem:A.1} and~\ref{lem:A.2} prove that the moment
matrices are invertible and that these finite systems have solutions. Each
time those lemmas are used, the needed bounds for the moment differences and
the matrix inverses are checked.
\end{EESegment}

\begin{EESegment}{EE-TGT-P030-008}
The moment conditions at infinity do a different job. The quantities
\NSi{NS-F-P030-031}{M,J,S} control, in that order, the radially integrated
axial momentum, the axial transport of angular momentum, and the axial
momentum flux including pressure. The angular moment
\NSi{NS-F-P030-032}{I} is made convergent by subtracting the prescribed
exterior power \NSi{NS-F-P030-033}{H_{\mathrm{pow}}} before integrating, as in
\eqref{eq:4.28}. The resulting total identities make both weighted integrals
of the tangential residuals vanish. Without those identities, the exterior
residual can be zero while the angular and axial stress components still have
tails proportional to \NSi{NS-F-P030-034}{r^{-2}} and
\NSi{NS-F-P030-035}{r^{-1}}, respectively, because the integrated inner
residual may still be nonzero. Lemma~\ref{lem:A.8} uses the total moment
identities to remove these tails and prove that the stress vanishes outside.
\end{EESegment}

\begin{EESegment}{EE-TGT-P030-009}
\subsection{The cone condition for the leading stress}\label{sec:4.3}

The annular stress \NSi{NS-F-P030-036}{\mathcal T_0} has to be produced by
the chosen waves with real amplitudes, so it must lie in the positive
span of their covariance directions, as Proposition~\ref{prop:7.5} proves.
The next formulas put this cone condition into profile variables, and
Lemma~\ref{lem:4.5} gives equivalent inequalities. At each
\NSi{NS-F-P030-037}{(X,\eta)}, the data that matter are the radial shear
\NSi{NS-F-P030-038}{(a,-b_s)} and the integrated inviscid term
\NSi{NS-F-P030-039}{p_s=(p_{s,1},p_{s,2})} from \eqref{eq:4.11}. Wherever
\NSi{NS-F-P030-040}{a>0}, set
\end{EESegment}
\NSFormula{NS-F-P030-041}{display}{4.20}
\begin{equation}
 t_s=-\frac{b_s}{a},\qquad
 v_s=a(1+t_s^2),\qquad
 P_c=p_{s,1}+t_sp_{s,2},\qquad
 J_c=p_{s,2}-t_sp_{s,1}.
\tag{4.20}\label{eq:4.20}
\end{equation}
